No aparato da Figura 3.12, o torque sobre o eixo de rotação é constante porque todos os parâmetros que o determinam são fixos.

\textbf{Fórmula do torque:}
Se uma massa $m$ está posicionada a uma distância radial $R$ fixada do eixo, o torque gravitacional é:
\[
\tau = m\,g\,R.
\]

\textbf{Por que é constante:}
\begin{itemize}
\item $m$: massa do elemento é fixada (não há perda ou ganho de massa);
\item $g$: aceleração gravitacional local constante;
\item $R$: braço de momento fixado rigidamente pelo desenho do aparato.
\end{itemize}

Assim, $\tau = mgR = \text{constante}$. Essa constante é a condição que permite integrar a equação do momento angular:
\[
\tau = I\,\frac{d\omega}{dt} = \text{constante} \implies \omega(t) = \frac{\tau}{I}\,t + \omega_0,
\]
mostando que a velocidade angular cresce linearmente com o tempo enquanto o torque é constante e não há atrito.

